% !TEX root = ../notes.tex

\documentclass[../notes.tex]{subfiles}

\begin{document}

\section{March 9}
Today we continue with our applications of slopes.
\begin{remark}
	Fix an odd prime $p$. The element $v_1$ (constructed on the homework) corresponds to a map $v_1\colon\Sigma^{2p-2}\mathbb S/p\to\mathbb S/p$. For example, we see that $(\mathbb S/p)_*(X)$ is an $\mathbb F_p[v_1]$-module for any $X$.
\end{remark}

\subsection{Morava \texorpdfstring{$K$}{ K}-Theory}
Let's start by explaining some more general quotients from last class.
\begin{proposition}
	Suppose that $R$ is an $\mathbb E_2$-ring for which $\pi_*R$ is supported in even degrees. For any $x_1,\ldots,x_n\in R$, the quotient $R/(x_1,\ldots,x_n)$ can be given the structure of an $\mathbb E_1$-$R$-algebra.
\end{proposition}
\begin{proof}
	For $R=\mathrm{MUP}$, we already saw how to do this for $\pi_2$ in \Cref{rem:quotient-by-pi2}. The same argument works for any $R$ for our elements in $\pi_2$.

	In general, there is a product decomposition $\mathrm{BU}=\CP^\infty\times\mathrm{BSU}$. (Thinking of $\mathrm{BU}$ as virtual vector bundles, one has a leftward map by taking actual line bundles, and one has a map from the right to the left inverting it.) Applying $\Omega^2$, we see that
	\[\Omega^2\mathrm{BU}=\ZZ\times\mathrm{BU}\]
	as $2$-fold loop spaces (by Bott periodicity). Now, one can take the Thom spectrum of the functor $\NN\stackrel m\to\ZZ\times\mathrm{BU}\to\op{Pic}\mathbb S$ to produce $\mathbb E_2$-rings $\mathbb S[x_{2m}]$. One then needs to solve some lifting property as in \Cref{rem:quotient-by-pi2}, which eventually follows from the fact that the even-ness of $R$ and the fact that $\CP^\infty$ has a cell structure using only cells of even dimension.
\end{proof}
\begin{remark}
	The lifting described at the end of the argument implies that the $\mathbb E_1$-structure is far from canonical.
\end{remark}
Here is our application.
\begin{definition}
	Recall that $\pi_*\mathrm{BP}=\ZZ_{(p)}[v_1,v_2,\ldots]$, where $v_i$ lives in degree $2p^i-2$. For each $n\ge0$, we define $\mathrm{BP}\langle n\rangle$ as the quotient
	\[\frac{\mathrm{BP}}{(v_{n+1},v_{n+2},\ldots)}.\]
\end{definition}
\begin{remark}
	By construction, we see that $\mathrm{BP}\langle n\rangle$ admits the structure of an $\mathbb E_1$-$\mathrm{BP}$-algebra, though it does not even admit a unique structure as a $\mathrm{BP}$-module because that requires a choice of generators $v_\bullet$s!
\end{remark}
\begin{remark}
	For $p>2$ or when $p=2$ and $n\le2$, it turns out that $\mathrm{BP}\langle n\rangle$ is unique as a spectrum. In these cases, it is reduced to some $\op{Ext}$ calculation, which is still open at $p=2$.
\end{remark}
We can further pass to ``residue fields.''
\begin{definition}
	We further define $k(n)\coloneqq\mathrm{BP}\langle n\rangle/(p,v_1,\ldots,v_{n-1})$ and $K(n)\coloneqq v_n^{-1}\mathrm{BP}\otimes_{\mathrm{BP}}k(n)$.
\end{definition}
\begin{remark}
	Note that $\pi_*k(n)\cong\FF_p[v_n]$, where $v_n$ lives in degree $2p^n-2$. It is known that $k(n)$ is unique as a spectrum.
\end{remark}
\begin{remark}
	By construction, we see that $\pi_*K(n)=\FF_p\left[v_n,1/v_n\right]$, so $K(n)$ is a skew field. In particular, every $K(n)$-module is free.
\end{remark}
\begin{remark}
	The $K(n)$s are called Morava $K$-theory. It turns out that $K(1)$ appears in the decomposition
	\[\mathrm{KU}/p=K(1)\oplus\Sigma^2K(1)\oplus\cdots\oplus\Sigma^{p-1}K(1),\]
	where we will choose not to define $\mathrm{KU}$.
\end{remark}

\subsection{Characteristics of Skew Fields}
We are now ready to classify characteristics of skew fields. To start, we note that we don't need to remove the lower $v$s to get a free module.
\begin{proposition} \label{prop:many-mods-is-still-free}
	The spectrum
	\[\left(\frac{\mathrm{BP}}{(p,v_1,\ldots,v_{n-1})}\right)[1/v_n]\]
	is an infinite direct sum of shifts of $K(n)$.
\end{proposition}
\begin{proof}
	Recall that there is a map $\eta_R\colon\mathrm{BP}\to\mathrm{BP}\otimes\mathrm{BP}$. On the other hand, $\pi_*(\mathrm{BP}\otimes\mathrm{BP})$ was a polynomial ring over $\pi_*\mathrm{BP}$, which we saw from a decomposition
	\[\mathrm{BP}\otimes\mathrm{BP}=\bigoplus_{t_I}\Sigma^{\deg t_I}\mathrm{BP}.\]
	Thus, we have a composite
	\[\mathrm{BP}\to\mathrm{BP}\otimes\mathrm{BP}=\bigoplus_{t_I}\Sigma^{\deg t_I}\mathrm{BP}\onto\bigoplus_{\text{monomials }t_I\in\pi_*\mathrm{BP}\left[t_1^{p^n},t_2^{p^n},\ldots\right]}\Sigma^{\deg t_I}\mathrm{BP}\langle n\rangle,\]
	where the last map is a projection. Modding out by $(p,v_1,\ldots,v_n)$ and then inverting $v_n$ produces a map from $\left(\frac{\mathrm{BP}}{(p,v_1,\ldots,v_{n-1})}\right)[1/v_n]$ to a sum of shifts of $K(n)$. One can then compute this composite on homotopy groups, where it is seen to be an isomorphism.
\end{proof}
Here is a useful notion for classification.
\begin{definition}[Bousfield class]
	Two spectra $E$ and $F$ are in the same \textit{Bousfield class} if and only if $E\otimes X=0$ is equivalent to $F\otimes X=0$ for all spectra $X$.
\end{definition}
\begin{proposition} \label{prop:bp-bousfield}
	The spectrum $\mathrm{BP}$ is in the same Bousfield class as
	\[\FF_p\oplus\QQ\oplus K(1)\oplus K(2)\oplus\cdots.\]
\end{proposition}
\begin{proof}
	We will use the following claim repeatedly: for any spectrum $E$ with endomorphism $\nu\colon\Sigma^kE\to E$, one has that $E$ and $E/\nu\oplus E[1/\nu]$ are in the same Bousfield class. Indeed, to see why this completes the proof, note that $\mathrm{BP}\otimes X$ vanishes if and only if both $\mathrm{BP}[1/p]\otimes X$ and $\mathrm{BP}/p\otimes X$ vanishes. The former is equivalent to $\QQ\otimes X$. Then the latter is equivalent to having $\mathrm{BP}/p[1/v_1]\otimes X=0$ and $\mathrm{BP}/(p,v_1)\otimes X=0$. The former is equivalent to $X\otimes K(1)=0$, and for the latter we keep going to $K(2)$, then $K(3)$, and so on. The $\FF_p$ appears because one should show that $X\otimes\FF_p$ vanishes if and only if $X\otimes K(n)$ vanishes for $n$ large enough. Indeed,
	\[\FF_p=\frac{\mathrm{BP}}{(p,v_1,v_2,\ldots)},\]
	which is some filtered colimit, so having $X\otimes\FF_p$ vanish means that it will vanish after some finite colimit, meaning that $X\otimes K(n)$ vanishes for $n$ large enough.

	It remains to show the claim. Certainly $E\otimes X$ vanishing implies that $E/\nu\otimes X$ and $E[1/\nu]\otimes X$ vanishes. Conversely, $E/\nu\otimes X$ vanishing means that the map $\nu\colon\Sigma^kE\otimes X\to E\otimes X$ is an equivalence, so $\nu$ acts by an isomorphism on $\pi_*(E\otimes X)$. But $E[1/\nu]\otimes X$ vanishes, so this isomorphism must vanish, so $\pi_*(E\otimes X)=0$, so $E\otimes X=0$.
\end{proof}
\begin{remark} \label{rem:check-ring-vanish-morava}
	Thus, by \Cref{ex:test-ring-zero-bp}, we see that a $p$-local $\mathbb E_1$-ring $R$ vanishes if and only if $R\otimes\FF_p=0$ or $R\otimes\QQ=0$ or $R\otimes K(n)=0$ for all $n\ge1$. Presumably, it is easier to check if tensoring with one of the fields $\{\FF_p,\QQ\}\cup\{K(n)\}_{n\ge1}$ vanishes is easier because one can use K\"unneth formulae and the like.
\end{remark}
We are now ready to classify some fields.
\begin{definition}[characteristic]
	Two skew fields $R_1$ and $R_2$ are in the same \textit{characteristic} if and only if $R_1\otimes R_2\ne0$.
\end{definition}
One can check that this is an equivalence relation, though it requires a little effort.
\begin{remark}
	It turns out that the skew fields $\{\QQ,\FF_p\}\cup\{K(n)\}_{n\ge1}$ all live in different characteristic classes.
\end{remark}
\begin{proposition}
	Every $p$-local skew field has the characteristic of $\QQ$, $\FF_p$, or $K(n)$ for some $n\ge1$.
\end{proposition}
\begin{proof}
	Let $R$ be a $p$-local skew field. Then $R$ is nonzero, so \Cref{rem:check-ring-vanish-morava} tells us that some element in
	\[\{R\otimes\QQ,R\otimes\FF_p\}\cup\{R\otimes K(n)\}_{n\ge1}\]
	is nonzero, so the statement follows.
\end{proof}
\begin{remark}
	One can remove the $p$-local hypothesis by simply letting $p$ vary.
\end{remark}

\subsection{Finite Spectra}
Let's next talk about finite spectra.
\begin{definition}[finite]
	A \textit{finite spectrum} is a perfect $\mathbb S_{(p)}$-module.
\end{definition}
Here are some quick facts.
\begin{lemma} \label{lem:finite-nonzero-by-fp}
	Fix a finite spectrum $X$. Then $X$ vanishes if and only if $X\otimes\FF_p=0$.
\end{lemma}
\begin{proof}
	For abelian groups, this is basically Nakayama's lemma. One then combines this with the Hurewisz theorem in order to show that $X\otimes\FF_p=0$ implies that we can kill the homotopy groups one at a time.
\end{proof}
\begin{lemma}
	Fix a nonzero finite spectrum $X$. Then $K(n)_*X\ne0$ for $n$ sufficiently large.
\end{lemma}
\begin{proof}
	There is an Atiyah--Hirzebruch spectral sequence
	\[\mathrm H_*(X;\FF_p[v_n,1/v_n])\Rightarrow K(n)_*X.\]
	Now, because $X$ is finite, $\pi_*(X\otimes\FF_p)=\mathrm H_*(X;\FF_p)$ is concentrated in only finitely many degrees (and is nonzero by \Cref{lem:finite-nonzero-by-fp}). The left-hand side is some sum of $\Sigma^{m\left(2p^n-2\right)}\mathrm H_*(X;\FF_p)$s, so there is a finite band in the plane of the spectral sequence in which the $E_2$-page concentrates. But as $n$ goes large, the differentials will be unable to be nonzero, so the spectral sequence immediately converges, and we see that $K(n)_*X$ is nonzero.
\end{proof}
\begin{lemma}
	Fix a finite spectrum $X$. If $K(n)\otimes X=0$, then $K(n-1)\otimes X=0$.
\end{lemma}
\begin{proof}
	Consider the $\mathbb E_1$-ring
	\[R=\left(\frac{\mathrm{BP}}{(p,v_1,\ldots,v_{n-2},v_{n+1},v_{n+2},\ldots)}\right)[1/v_n]\]
	so that $\pi_*R=\FF_p[v_{n-1},v_n,1/v_n]$. Because $X$ is perfect, $\pi_*(R\otimes X)$ is a finitely generated module over $\pi_*R$, and it vanishes$\pmod{v_{n-1}}$. But $\pi_*R$ is some kind of polynomial ring in one variable over a graded field, so such (graded) modules can be classified. Using this classification, one is able to see that vanishing$\pmod{v_{n-1}}$ implies that $\pi_*(R\otimes X)[1/v_{n-1}]$ also vanishes. Indeed, $R_*X$ is $\left(2p^n-2\right)$-periodic (because $v_n$ has invertible action), and it is a module over $\pi_0R$, which is $\FF_p[t]$ where $t=v_{-1}^dv_n^{-1}$ for some $d$. As such, we are looking at $2p^n-2$ modules over $\FF_p[t]$, and our conditions that we vanish$\pmod{v_{n-1}}$ and vanish after inverting $v_{n-1}$ gather together to show that our module over a principal ideal domain implies that $K(n-1)\otimes X=0$.
\end{proof}
Combining these lemmas grants us an invariant.
\begin{definition}[type]
	Fix a finite spectrum $X$. Then the \textit{type} of $X$ is the unique nonnegative integer $n$ for which $K(m)\otimes X=0$ for $m<n$ and $K(m)\otimes X\ne0$ for $m\ge n$. By convention, $K(0)=\QQ$.
\end{definition}
\begin{example}
	One has that $\mathbb S_{(p)}$ is type $0$ because $\mathbb S_{(p)}\otimes\QQ\ne0$.
\end{example}
\begin{example}
	Note that $\mathbb S/p$ is type $1$. It is not hard to see that $\QQ\otimes\mathbb S/p=0$, and one can check that $K(1)\otimes\mathbb S/p\ne0$.
\end{example}
\begin{example}
	For odd primes $p$, we showed on the homework that the finite spectrum $\mathbb S/(p,v_1)$ is type $2$.
\end{example}
\begin{example} \label{ex:s-eta-type-zero}
	Let $\eta$ be the Hopf fibration. Then $\mathbb S/\eta$ is type zero: $\eta$ is a map $\Sigma\mathbb S\to\mathbb S$, so $\QQ\otimes\mathbb S/\eta$ is the cofiber of some map $\Sigma\QQ\to\QQ$. But all such maps vanish, so $\QQ\otimes\mathbb S/\eta=\QQ\oplus\Sigma^2\QQ$, which is nonzero.
\end{example}
\begin{remark}
	It was verified (by brute force construction) that there is a finite spectrum of each type. One basically takes finite retracts of tensor products of some well-understood ``lens spaces.''
\end{remark}

\end{document}